\begin{description}
\item[(a)] Since $E \propto M \rightarrow L\tau \propto M$ and
$L \propto M^{3.5}$,
\[ \tau = \tau_\odot\left(\frac{M_\odot}{M}\right)^{2.5} \]
\hfill(2.0)

The maximum mass of the star should be
\[ M = M_\odot\left(\frac{\tau_\odot}{\tau}\right)^{1/2.5}
     = M_\odot\left(\frac{10^{10}}{4\times10^{9}}\right)^{1/2.5} \]
\hfill(0.5)
\[ M = 1.44 M_\odot \]
\hfill(0.5)

Use of spectral class--$\log M$ linear relation
\[ \frac{10}{(\log 1.6 - \log 1.05)}\times(\log 1.6 - \log 1.44) = 2.5 \]
\hfill(2.0)

which corresponds to the star of type \textbf{F2V} to \textbf{F3V}. \hfill(1.0)

\textit{0.5 mark for the answer ``between F0V and G0V''.}

\item[(b)] The received power by a planet of a radius $r$ and a distance $d$
from the star is
\[ P_{input} = (1-a)\frac{\pi r^2 L}{4\pi d^2} = (1-a)\frac{r^2 L}{4d^2}, \]
where $a$ is the planet's albedo. \hfill(1.0)

The released power of blackbody radiation is
\[ P_{output} = 4\pi r^2\epsilon\sigma T^4, \]
where $\epsilon$ is emissivity. \hfill(1.0)

Since we assume that the planet emits heat as blackbody radiation,
\[ (1-a)\frac{r^2 L}{4d^2} = 4\pi r^2\epsilon\sigma T^4. \]
\hfill(1.0)

The temperatures of the planet and the Earth are similar,
\[ T^4 = \frac{(1-a)L}{16\pi\epsilon\sigma d^2}
       = \frac{(1-a)L_\odot}{16\pi\epsilon\sigma d_\odot^2}. \]
\[ \frac{L}{d^2} = \frac{L_\odot}{d_\odot^2} \]
\hfill(1.0)
\[ \left(\frac{d}{d_\odot}\right)^2 = \frac{L}{L_\odot}
   = \left(\frac{M}{M_\odot}\right)^{3.5} \]
\hfill(1.0)

The distance of the planet to the star is
$d = \left(\frac{M}{M_\odot}\right)^{3.5/2}$~au. \hfill(1.0)

\item[(c)] Angular speed around the center of mass can be derived by Kepler's
law.
\[ T^2 = \frac{4\pi^2 d^3}{G(M+m)} \]
\hfill(1.0)
\[ \omega = \sqrt{\frac{G(M+m)}{d^3}} \]
\hfill(1.0)

According to the centre of mass of the planet and the star,
\begin{align*}
Md_s &= md\\
d_s &= \frac{m}{M}\,d
\end{align*}
\hfill(1.0)

and
\[ v = \omega d_s \]
\hfill(1.0)
\[ v = \omega d\,\frac{m}{M} = \frac{m}{M}\sqrt{\frac{G(M+m)}{d}}
     \approx m\sqrt{\frac{G}{Md}} \]

From
\[ \frac{\Delta\lambda}{\lambda} = \frac{v}{c} = \frac{m}{c}\sqrt{\frac{G}{Md}} \]
\hfill(1.0)

Use of previously derived expression. The distance of the planet to the star
is
\[ d = \left(\frac{M}{M_\odot}\right)^{3.5/2}\ \mathrm{au}. \]
\hfill(1.0)
\begin{align*}
m &= \frac{\Delta\lambda}{\lambda}\,c\sqrt{\frac{Md}{G}}
   = \frac{\Delta\lambda}{\lambda}\,c
     \sqrt{\frac{M\left(\frac{M}{M_\odot}\right)^{3.5/2}\times 1\ \mathrm{au}}{G}}\\
  &= \frac{\Delta\lambda}{\lambda}\,c
     \sqrt{\frac{M\,(1.44)^{3.5/2}\times 1\ \mathrm{au}}{G}}\\
  &= (299792458\ \mathrm{m/s})\times10^{-10}
     \sqrt{\frac{1.44\times1.99\times10^{30}\ \mathrm{kg}\,
       (1.44)^{3.5/2}\times1.50\times10^{11}\ \mathrm{m}}
       {6.67\times10^{-11}\ \mathrm{m^3kg^{-1}s^{-2}}}}
\end{align*}
\hfill(1.0)
\[ m = 3.31\times10^{24}\ \mathrm{kg} = 0.554 M_\oplus \]
\hfill(1.0)
\end{description}
